\section{Harmonic restriction with a continuous base metric}\label{har:sec:harmonic}

The coefficient on the total space in the decomposition theorem below
always has a smooth metric. Continuity enters only through a line bundle
pulled back from the base. We give the approximation and restriction
arguments because the decomposition theorem by itself is not a vanishing
theorem.

\begin{lemma}[Uniform smoothing]\label{har:lem:smoothing}
Let \(Q\) be a line bundle on a smooth projective \(Y\), with a positive
continuous Hermitian metric whose local weights are plurisubharmonic.
For a fixed Kähler form \(\eta\), there are smooth metrics whose weights
converge uniformly to these weights and whose curvatures are at least
\(-\varepsilon_j\eta\), with \(\varepsilon_j\to0\).
\end{lemma}

\begin{proof}
This is the regularized-maximum construction of Richberg
\cite{Richberg1968}; see also \cite[Theorem~I.5.21]{Demailly2012}. The small negative curvature allowance permits its
use for semipositive, rather than strictly positive, weights. Write a
weight as a smooth reference weight plus a global continuous function
\(u\), with background curvature \(\gamma\). For fixed \(\varepsilon>0\),
use background \(\gamma+\varepsilon\eta\), so that the original potential
has a strict local margin. On finitely many enlarged coordinate patches,
add a smooth local potential of this background, convolve, and subtract
that potential. Take inner patches which still cover \(Y\). Add smooth
bumps zero on the inner patches and a small negative constant near the
outer boundaries, choosing their second derivatives smaller than the
strict margin and the convolution errors much smaller than the bump
heights. A regularized maximum of the active candidates is smooth: a
candidate near its outer boundary is uniformly below a candidate from
an inner patch and does not affect the maximum. Convexity, monotonicity,
and translation equivariance of regularized maximum preserve the
background lower bound and the transition rule for weights. Let the
uniform error and then \(\varepsilon\) tend to zero. Equivalently, one
may smooth the strictly positive metric on \(NQ+H\), with \(H\) positive,
subtract the weight of \(H\), and divide by \(N\).
\end{proof}

\begin{theorem}[Kernel-condition vanishing]\label{har:thm:kernel}
Let $f:V\to Y$ be a surjective morphism between smooth connected projective complex varieties. Let $D_V$ be a line bundle with a smooth semipositive Hermitian metric, whose curvature is $\theta$. Let $Q$ be a line bundle on $Y$ with a positive continuous Hermitian metric having plurisubharmonic weights.
Assume that on a nonempty ordinary open subset of the submersion locus,
\begin{equation}\label{har:eq:kernel}
 \ker\theta\subseteq\ker df.
\end{equation}
Then, for every \(k\ge0\), \(q\ge0\), and \(p>0\),
\begin{equation}\label{har:eq:vanishing}
 H^p\bigl(Y,R^q f_*(K_V+D_V)\otimes Q^k\bigr)=0.
\end{equation}
In applications below \(D_V\) is the pullback of a smoothly semipositive
anticanonical line, and the particular diagram supplies
\eqref{har:eq:kernel}; the rank-open condition is not inferred from nefness.
\end{theorem}

\begin{proof}
Fix \(k\). We will show that every nonzero cohomology class on the
source restricts nontrivially to some general fiber. Smooth-coefficient
decomposition will then give the asserted vanishing. Choose a Kähler
form \(\omega\) on \(V\) and a Dolbeault class
\(c\in H^\ell(V,K_V+A)\), where \(A=D_V+kf^*Q\) and \(n=\dim V\).
Lemma~\ref{har:lem:smoothing} gives product metrics \(h_j\to h\) uniformly
on \(A\), with curvature at least \(\theta-o(1)\omega\). Let \(u_j\)
be the harmonic \((n,\ell)\)-representative of \(c\) for \(h_j\).
Harmonic norm minimization, comparison with one fixed smooth closed
representative, and uniform metric equivalence give a fixed \(L^2\) bound.
The Bochner--Kodaira identity gives
\begin{equation}\label{har:eq:energy}
 \|(D'_j)^*u_j\|_{h_j}^2
 +\langle[\theta,\Lambda]u_j,u_j\rangle_{h_j}
 \le o(1)\|u_j\|_{h_j}^2.
\end{equation}
Here curvature commutators in bidegree \((n,\ell)\) are monotone in the
curvature form: their eigenvalues are sums of the curvature eigenvalues
on the antiholomorphic index set.

Let \(v_j=*u_j\), absorbing a fixed nonzero dimensional constant so that
\(u_j=\omega^\ell\wedge v_j\). The complex-linear star acts on form
indices, carrying the coefficient line along. Metric compatibility gives
\((D'_j)^*=\pm*\bar\partial*\) and
\(\bar\partial_j^*=\pm*D'_j*\). Thus
\[
 \bar\partial v_j\longrightarrow0\text{ in }L^2,
 \qquad D'_jv_j=0.
\]
These are statements about a fixed holomorphic bundle; the first
operator does not differentiate the varying metric. Passing to a weak
\(L^2\) subsequence gives a distributionally holomorphic, hence smooth
holomorphic, \(v\). Since \(\omega\) is Kähler, the form
\(u=\omega^\ell\wedge v\) is \(\bar\partial\)-closed and represents
\(c\). Indeed projection to harmonic forms for one fixed smooth
reference metric is a bounded finite-rank operator, has the same value
on every \(u_j\), and commutes with this weak limit. The nonnegative
square root of \([\theta,\Lambda]\) annihilates the limit by \eqref{har:eq:energy}.
Pointwise diagonalization therefore shows that every positive-curvature
holomorphic covector is a wedge factor of \(v\).

The approximation also gives a Chern equation before any use of the
kernel condition. In any local holomorphic frame of \(A\), let
\(h_j\) denote its squared norm. The equation \(D'_jv_j=0\) reads
\(\partial(h_jv_j)=0\). Uniform convergence of \(h_j\) and weak
\(L^2\) convergence suffice to pass it to distributions:
\begin{equation}\label{har:eq:distribution}
 \partial(hv)=0.
\end{equation}
No derivative convergence of metric weights is used.

To detect the class on a fiber, we next show that \(v\) contains all
base canonical covectors as wedge factors. On the open set in
\eqref{har:eq:kernel}, every pulled-back base covector annihilates
\(\ker\theta\), so its wedge with \(v\) vanishes. The holomorphic
bundle map
\[
 f^*\Omega_Y^1\longrightarrow\Omega_V^{n-\ell+1}\otimes A,
 \qquad \alpha\longmapsto (df)^*\alpha\wedge v,
\]
therefore vanishes everywhere. Over the smooth fibration locus, in
base coordinates \(y_1,\ldots,y_b\), this gives
\begin{equation}\label{har:eq:basefactor}
 v=f^*(dy_1\wedge\cdots\wedge dy_b)\wedge w.
\end{equation}
The relative restriction \(w_y\) is unambiguously defined after choosing
the base canonical frame. If \(n-\ell<b\), then \(v=0\) already.

The limiting metric is the smooth metric of \(D_V\) times a positive
continuous function pulled back from the base. Because of
\eqref{har:eq:basefactor}, only vertical differentiations survive in
\eqref{har:eq:distribution}.
The continuous base function is constant under these differentiations;
the remaining expression is smooth in the vertical variables and
continuous altogether. It vanishes pointwise and gives
\(D'_{D_V|V_y}w_y=0\) on every smooth fiber under consideration.

The restriction in cohomology is restriction of a \((0,\ell)\)-form
\emph{valued in} \(K_V+A\). Use
\(K_V|_{V_y}=K_{V_y}\otimes K_Y|_y\) to remove the constant canonical
base factor. It is not the ordinary pullback of an \((n,\ell)\)-form to
a lower-dimensional fiber. The restricted representative is a nonzero
constant times
\[
 \omega_y^\ell\wedge w_y,
 \qquad \omega_y=\omega|_{V_y}.
\]
It is harmonic on the compact smooth fiber: it is \(\bar\partial\)-closed,
and its star has zero Chern derivative. Thus a nonzero \(v\) gives a
nonzero restriction class on some general fiber. The Leray edge map
\begin{equation}\label{har:eq:edge}
 H^\ell(V,K_V+D_V+kf^*Q)
 \longrightarrow H^0(Y,R^\ell f_*(K_V+D_V)\otimes Q^k)
\end{equation}
is injective, since zero edge image implies zero restriction on the
general smooth base-change locus.

Finally apply Fujisawa's derived decomposition
\cite[Theorem~1]{Fujisawa2015}, obtained from Takegoshi's harmonic
representatives, to the \emph{smooth} coefficient \(D_V\):
\[
 Rf_*(K_V+D_V)\simeq\bigoplus_q R^q f_*(K_V+D_V)[-q].
\]
Tensor the resulting derived isomorphism by the line \(Q^k\) on \(Y\)
and use the projection formula. It gives
\[
 h^\ell(V,K_V+D_V+kf^*Q)
 =\sum_{p+q=\ell}h^p(Y,R^qf_*(K_V+D_V)\otimes Q^k).
\]
Injection \eqref{har:eq:edge} bounds the left side by the summand \(p=0\).
All remaining summands vanish. For total degrees exceeding \(n\), the
left side is zero directly. This proves every case of
\eqref{har:eq:vanishing}, including the twist \(k=0\).
\end{proof}

\begin{remark}[The vertical-polyvector proof]\label{har:rem:polyvector}
There is a second way to obtain the restriction step, useful when only
uniform comparability, rather than convergence of metrics, is known.
Let \(\rho\) be a smooth semipositive real \((1,1)\)-form on \(V\). Suppose a line \(A\) has uniformly comparable smooth metrics of curvature
\(\ge c\rho-\varepsilon_j\omega\), where \(c>0\) is fixed and \(\varepsilon_j\to0\), on a smooth projective \(V\).
The weak-limit argument through \eqref{har:eq:energy} gives a holomorphic
\(v\in H^0(\Omega_V^{n-\ell}\otimes A)\) representing each adjoint class
by \(\omega^\ell\wedge v\). Under
\(\Omega_V^{n-\ell}\otimes A=\bigwedge^\ell T_V\otimes(K_V+A)\),
it corresponds to a holomorphic polyvector whose indices lie in
\(\bigwedge^\ell\ker\rho\). If \(\ker\rho\subset\ker df\) on a
nonempty open, this polyvector is vertical throughout the submersion
locus by holomorphy. Lowering its vertical indices and restricting uses
exactly the induced fiber Kähler metric, with the canonical base factor
held constant. Suppose on each general fiber there is a smooth metric
on \(A|_{V_y}\) with curvature a positive multiple of \(\rho|_{V_y}\).
The corresponding holomorphic \((\dim V_y-\ell,0)\)-form \(v_y\)
contains all the positive-curvature covectors. Hence its curvature
commutator is zero. Since a holomorphic form of antiholomorphic degree
zero is \(\bar\partial\)-harmonic, Bochner--Kodaira gives \(D'v_y=0\).
Its Lefschetz image is fiber-harmonic. This proves the same detection
and edge injection without passing a Chern equation through the limit.
The smooth-coefficient splitting is still a separate requirement.
\end{remark}

\begin{proposition}[Pseudoeffective domination variant]\label{har:prop:dominationvanish}
In Theorem~\ref{har:thm:kernel}, replace the kernel hypothesis by the
assumption that \(bD_V-f^*H\) is pseudoeffective for some \(b>0\) and
some ample line \(H\) on \(Y\). Then the same vanishing holds.
\end{proposition}

\begin{proof}
Keep the uniformly convergent metrics and the notation of the proof of
Theorem~\ref{har:thm:kernel}. For \(\ell>0\), put
\[
 S=i^{(n-\ell)^2}h\,v\wedge\overline v,
 \qquad T=S\wedge\omega^{\ell-1},
\]
with coefficient pairing understood. The distributional Chern equation
and holomorphy make \(T\) a continuous closed positive
\((n-1,n-1)\)-form. The curvature-energy equality gives
\(\theta\wedge T=0\). Choose a positive curvature form \(\eta\) of
\(H\). A positive current in the class of
\(bD_V-f^*H\) pairs nonnegatively with \(T\), since positive currents
have order zero. To compute this pairing cohomologically, apply de Rham
regularization, which commutes with \(d\), to approximate \(T\)
uniformly by smooth closed forms. For each approximant the pairing
agrees with that against the cohomologous smooth form
\(b\theta-f^*\eta\); uniform convergence passes this equality to
\(T\). It follows that
\[
 0\le\int_V(b\theta-f^*\eta)\wedge T
       =-\int_V f^*\eta\wedge T\le0.
\]
Pointwise positivity forces \(f^*\eta\wedge T=0\); diagonalization gives
the full base wedge factor in \(v\). The distributional fiber equation,
canonical-valued restriction, edge injection, and smooth-coefficient
splitting proceed exactly as in that proof. Only positive total degrees
are needed to kill a term of positive base degree.
\end{proof}

\begin{proposition}[Euler interpolation on a birational model]\label{har:prop:euler}
Let \(\pi:V\to X\) be birational between smooth connected projective
varieties, let \(L=-K_X\) be smoothly semipositive, put
\(D_V=\pi^*L\), and write \(E=K_{V/X}=K_V+D_V\).
Let \(f:V\to Y\) be surjective, with \(Y\) smooth connected
projective. Suppose a continuously semipositive line \(Q\) on \(Y\)
and an injection \(f^*Q\to bD_V\), \(b>0\), satisfy either vanishing
criterion above. Then \(H^0(X,kbL)\ne0\) for some \(k>0\).
If a torus acts, trivially on \(Y\) and \(Q\), and all maps and the
injection are equivariant, this section can be taken invariant.
\end{proposition}

\begin{proof}
Let \(\mathcal G=f_*\mathcal O_V(E)\), or its invariant summand in the
equivariant case. A torus over a trivially acted-on base has coherent
weight summands, so this is a direct summand. Vanishing gives
\(h^0(Y,\mathcal G\otimes Q^k)=\chi(Y,\mathcal G\otimes Q^k)\) for every
\(k\ge0\). The right side is polynomial in \(k\) by Riemann--Roch for
coherent sheaves. At zero it is one: \(\pi_*\mathcal O_V(E)=\mathcal O_X\)
and the exceptional canonical section is invariant. The polynomial is
therefore nonzero at some positive integer. Multiply the resulting
section of \(E+kf^*Q\) by the \(k\)-th power of the injection and push
through \(\pi\). The effective exceptional divisor \(E\) changes neither
sections nor their characters, by normality of \(X\).
\end{proof}
